%% kapa: exam=1 round=33 year=2022 period=1 form=A subject=economics subject_ko=경제학원론 seq=3 number=43
%% source: Q-Net 감정평가사 기출문제 (cst00309) artlSeq=5209997 file=raw/1008/5209997/2022년도 제33회 감정평가사 1차시험 1교시(민법,경제학,부동산학).pdf
%% conversion: kordoc 4.24.51 -> tools/gen1.py (auto); review=proofread-sonnet; flags=none
\begin{question}{43}
( )에 들어갈 내용으로 옳은 것은?
\begin{framedtext}
○ 소비의 긍정적 외부성이 존재할 때, ( ㄱ )이 ( ㄴ )보다 크다. \par
○ 생산의 부정적 외부성이 존재할 때, ( ㄷ )이 ( ㄹ )보다 작다.
\end{framedtext}
\begin{choices}
  \item ㄱ: 사회적 한계편익, ㄴ: 사적 한계편익, ㄷ: 사적 한계비용, ㄹ: 사회적 한계비용
  \item ㄱ: 사적 한계편익, ㄴ: 사회적 한계편익, ㄷ: 사적 한계비용, ㄹ: 사회적 한계비용
  \item ㄱ: 사회적 한계편익, ㄴ: 사적 한계편익, ㄷ: 사회적 한계비용, ㄹ: 사적 한계비용
  \item ㄱ: 사적 한계편익, ㄴ: 사회적 한계편익, ㄷ: 사회적 한계비용, ㄹ: 사적 한계비용
  \item ㄱ: 사회적 한계편익, ㄴ: 사적 한계비용, ㄷ: 사적 한계편익, ㄹ: 사회적 한계비용
\end{choices}
\end{question}
\answer{1}
